%=========== ΛΥΣΗ - 55ο ΘΕΜΑ Γ ===========
\begin{thema}{Γ}
\begin{erwthma}
\item Θεωρούμε τη συνάρτηση
\[f(x)=e^x+x-2\ ,\ x\in\mathbb{R}\]
Η $f$ είναι συνεχής στο $[0,1]$ ως άθροισμα συνεχών συναρτήσεων.
Έχουμε
\[f(0)=e^0+0-2=-1<0\]
και
\[f(1)=e+1-2=e-1>0\]
Άρα
\[f(0)f(1)<0\]
Επομένως, από το θεώρημα \eng{Bolzano}, υπάρχει τουλάχιστον μία ρίζα $x_0\in(0,1)$ τέτοια ώστε
\[f(x_0)=0\]
δηλαδή
\[e^{x_0}+x_0=2\]

\item Η $f$ είναι παραγωγίσιμη στο $\mathbb{R}$ με
\[f'(x)=e^x+1>0\ ,\ x\in\mathbb{R}\]
Άρα η $f$ είναι \textbf{γνησίως αύξουσα} στο $\mathbb{R}$.
Επομένως η εξίσωση $f(x)=0$, δηλαδή η εξίσωση
\[e^x+x=2\]
έχει το πολύ μία ρίζα στο $\mathbb{R}$.
Αφού από το Γ1 γνωρίζουμε ότι έχει τουλάχιστον μία ρίζα $x_0\in(0,1)$, η ρίζα αυτή είναι μοναδική στο $\mathbb{R}$.

\item Από το προηγούμενο ερώτημα η $f$ είναι γνησίως αύξουσα στο $\mathbb{R}$ και
\[f(x_0)=0\]
Άρα
\begin{itemize}
\item για $x<x_0$ είναι $f(x)<f(x_0)=0$
\item για $x=x_0$ είναι $f(x)=0$
\item για $x>x_0$ είναι $f(x)>f(x_0)=0$
\end{itemize}
Επομένως
\[f(x)<0\ ,\ x\in(-\infty,x_0)\]
και
\[f(x)>0\ ,\ x\in(x_0,+\infty)\]

\item Η δοσμένη συνάρτηση είναι
\[F(x)=e^x+\frac{x^2}{2}-2x\ ,\ x\in\mathbb{R}\]
Παρατηρούμε ότι
\[F'(x)=e^x+x-2=f(x)\]
Από το Γ3 έχουμε
\begin{itemize}
\item $F'(x)<0$ για $x<x_0$
\item $F'(x)=0$ για $x=x_0$
\item $F'(x)>0$ για $x>x_0$
\end{itemize}
Άρα η $F$ είναι γνησίως φθίνουσα στο $(-\infty,x_0]$ και γνησίως αύξουσα στο $[x_0,+\infty)$.
Επομένως η $F$ παρουσιάζει ολικό ελάχιστο στη θέση $x_0$.

Επειδή
\[F'(x)=f(x)\]
έχουμε
\begin{align*}
\dint_0^{x_0} f(x)\d x&=F(x_0)-F(0)\\
&=\left(e^{x_0}+\frac{x_0^2}{2}-2x_0\right)-1
\end{align*}
Από τη σχέση $e^{x_0}+x_0=2$ προκύπτει
\[e^{x_0}=2-x_0\]
Άρα
\begin{align*}
\dint_0^{x_0} f(x)\d x
&=2-x_0+\frac{x_0^2}{2}-2x_0-1\\
&=\frac{x_0^2}{2}-3x_0+1
\end{align*}
Επομένως
\[\dint_0^{x_0} f(x)\d x=\frac{x_0^2}{2}-3x_0+1\]
\end{erwthma}
\end{thema}
%=========== ΤΕΛΟΣ ΛΥΣΗΣ - 55ο ΘΕΜΑ Γ ===========
